
Our experimental design is motivated by a single question: \textit{if a sampling-based consistency detector fails on an instruction-tuned model, is the failure in the scoring mechanism or in the underlying consistency assumption?} This requires a design that separates implementation correctness, scorer validity, and mechanism validity as independent diagnostic layers.

\subsubsection{Overview}

We evaluate NLI-based Semantic Mode Consistency (SMC-NLI), a direct instantiation of the SelfCheckGPT-NLI scoring protocol \citep{Manakul2023SelfCheckGPT}, applied to Llama-3-8B-Instruct on HaluEval QA \citep{Li2023HaluEval}. The core pipeline computes pairwise NLI agreement across N=10 stochastic samples per question and uses this consistency score as a hallucination predictor. We extend this with two design choices enabling clean attribution: a parallel embedding-based scorer (SMC-Embed) and an explicit mechanism verification step.

\subsubsection{Dataset}

\textbf{HaluEval QA} \citep{Li2023HaluEval} consists of question-answer pairs with binary hallucination labels (correct / hallucinated), constructed by generating ChatGPT responses and using ChatGPT to create hallucinated alternatives. We use a stratified sample of $N=1000$ questions (500 correct, 500 hallucinated), drawn with \texttt{seed=42}, providing adequate statistical power for AUROC estimation (standard error < 0.02 at this sample size).

HaluEval is among the most favorable structured factual QA benchmarks for this evaluation --- balanced labels, short-answer format, no annotation noise from human raters. We note the benchmark validity concern (ChatGPT labels vs. Llama-3-8B behavior) in Section 6.

\subsubsection{LLM Sampling}

For each question $q_i$, we generate $N=10$ stochastic samples from Llama-3-8B-Instruct (\texttt{meta-llama/Meta-Llama-3-8B-Instruct}) [Meta AI, 2024]:

$$S_i = \{s_i^{(1)}, s_i^{(2)}, \ldots, s_i^{(10)}\}$$

using temperature $\tau = 0.7$, top-p $p = 0.9$, and \texttt{max\textbackslash \_new\textbackslash \_tokens=50}. The model is loaded via HuggingFace Transformers with \texttt{device\textbackslash \_map=''auto''} for GPU memory management.

Temperature=0.7 matches the setting used in SelfCheckGPT \citep{Manakul2023SelfCheckGPT}, enabling methodological comparability. N=10 follows Kuhn et al. [2023], who found this sufficient for reliable semantic entropy estimation. Intermediate samples are saved to disk (\texttt{data/llama\textbackslash \_samples.json}) with resume logic to avoid re-inference.

Total inference: 10,000 samples across 1,000 questions.

\subsubsection{SMC-NLI Scoring}

For each question $q_i$ and its samples $S_i$, we compute the Semantic Mode Consistency (NLI variant) as the mean pairwise non-contradiction rate over all $\binom{N}{2} = 45$ sample pairs:

$$\text{SMC-NLI}(q_i) = \frac{1}{45} \sum_{j < k} \left[ P_\text{ent}(s_i^{(j)}, s_i^{(k)}) + P_\text{neut}(s_i^{(j)}, s_i^{(k)}) \right]$$

where $P_\text{ent}$ and $P_\text{neut}$ are entailment and neutral probabilities from the DeBERTa-v3-large cross-encoder NLI model (\texttt{cross-encoder/nli-deberta-v3-large}). Higher SMC-NLI indicates greater semantic consistency; lower SMC-NLI indicates diversity (inconsistency).

To mitigate the documented out-of-distribution concern for NLI models on short factual answers (trained on paragraph-length SNLI/MultiNLI pairs), each sample is prepended with the question context: \texttt{''Q: \{question\} A: \{answer\}''}, following Manakul et al. [2023].

NLI pairs are processed in batches of 16 on GPU. Total NLI pairs scored: 45,000.

\textbf{Prediction direction:} Higher SMC-NLI predicts correctness (label=0); lower SMC-NLI predicts hallucination (label=1). AUROC is computed in this direction.

\subsubsection{SMC-Embed Scoring (Parallel Control)}

In parallel with SMC-NLI, we compute an embedding-based consistency score:

$$\text{SMC-Embed}(q_i) = \frac{1}{45} \sum_{j < k} \cos\left(\mathbf{e}_i^{(j)}, \mathbf{e}_i^{(k)}\right)$$

where $\mathbf{e}_i^{(j)}$ is the sentence embedding of sample $s_i^{(j)}$ from \texttt{sentence-transformers/all-mpnet-base-v2}, and $\cos(\cdot, \cdot)$ denotes cosine similarity.

SMC-Embed measures the same underlying construct (output consistency) without the NLI scorer. If SMC-NLI fails due to NLI OOD for short factual answers, SMC-Embed should achieve meaningfully higher AUROC. If both fail at similar levels, the failure is in the shared assumption --- that consistent outputs carry no hallucination signal --- not in the NLI scorer specifically. This dual-metric design enables clean causal attribution of the negative result.

\subsubsection{Mechanism Verification}

Before full-scale evaluation, we perform a mechanism verification step on the first 5 questions:

\begin{enumerate}
\item Compute SMC-NLI scores for questions $\{q_1, \ldots, q_5\}$
\item Assert: all scores in $[0, 1]$
\item Assert: $\max_i \text{SMC-NLI}(q_i) - \min_i \text{SMC-NLI}(q_i) > 0.01$ (meaningful variation exists)
\item Raise an exception if either assertion fails
\end{enumerate}

This step distinguishes two failure types: ''the scorer is broken'' (no variation, all scores identical or out-of-range) vs. ''the scorer works but the scores are not discriminative at scale.'' When the mechanism check passes but full-scale AUROC collapses to chance, the failure is definitively in the hypothesis, not the implementation.

\subsubsection{Evaluation}

AUROC for both SMC-NLI and SMC-Embed is computed against HaluEval binary labels using \texttt{sklearn.metrics.roc\textbackslash \_auc\textbackslash \_score}. We additionally report:
\begin{itemize}
\item SMC-NLI standard deviation across questions (>0.05 indicates meaningful variation)
\item Per-label mean SMC-NLI scores (expected to differ between correct and hallucinated if the stochastic hallucination assumption holds)
\item ROC curves for both metrics
\item Per-question SMC-NLI vs. SMC-Embed scatter plot (correlation analysis)
\end{itemize}

\subsubsection{Implementation Validation}

All components are validated by 14 unit tests covering:
\begin{itemize}
\item Data pipeline correctness (stratified sampling, label extraction)
\item LLMSampler output format and seed reproducibility
\item SMCNLIScorer score range and batch consistency
\item SMCEmbedScorer cosine similarity range
\item Evaluate module AUROC computation
\end{itemize}

The Coder-Validator cycle completed in 1/5 iterations (all checks passing on first validation pass), providing independent confirmation that implementation errors are not a confound for interpreting the AUROC results.

\subsubsection{Summary of Key Design Choices}

\begin{table}[htbp]
\centering
\resizebox{\columnwidth}{!}{%
\begin{tabular}{ll}
\toprule
Design Choice & Rationale \\
\midrule
Dual metric (NLI + Embed) & Rules out NLI OOD as confound; enables regime-level attribution \\
Mechanism verification (5 questions) & Distinguishes broken code from wrong hypothesis \\
N=1000 balanced questions & Adequate statistical power (AUROC SE < 0.02) \\
14 unit tests & Independent implementation correctness verification \\
Question-prepend for NLI & Mitigates NLI OOD for short factual answers \\
Resume logic for samples & Reproducibility: exact sample set fixed, scorer can be rerun \\
\bottomrule
\end{tabular}
}
\end{table}

